\documentclass[12pt]{article}
\usepackage[utf8]{inputenc}
\usepackage[T1]{fontenc}
\usepackage[margin=1in]{geometry}
\usepackage{setspace}
\usepackage{xcolor}
\usepackage[colorlinks=true, urlcolor=blue!60!black, linkcolor=blue!60!black]{hyperref}
\setlength{\parindent}{2em}
\setlength{\parskip}{1.25em}
\renewcommand{\baselinestretch}{1.0}
\begin{document}
\begin{center}
{\LARGE \textit{Wisconsin v. Mitchell}, 508 U.S. 476 (1993)\par}
\vspace{0.8em}
{\large Case Brief\par}
\end{center}
\vspace{1em}

\section{Memory Jogger}
Unlike laws that criminalize bias-motivated speech, a hate-crime penalty-enhancement statute that increases the punishment for unprotected criminal conduct when the defendant selects the victim because of race or another protected characteristic does not violate the First Amendment.

\section{Detailed Case Facts}
On the evening of October 7, 1989, Todd Mitchell, a young black man, gathered with a group of young black men and boys at an apartment complex in Kenosha, Wisconsin (\href{https://supreme.justia.com/cases/federal/us/508/476/}{Justia Opinion}; \href{https://www.oyez.org/cases/1992/92-515}{Oyez}). Several members of the group discussed a scene from the motion picture \textit{Mississippi Burning}, in which a white man beat a young black boy who was praying. After the group moved outside, Mitchell asked, ``Do you all feel hyped up to move on some white people?'' Shortly thereafter, a young white boy walked past the group on the opposite side of the street. Mitchell said, ``You all want to fuck somebody up? There goes a white boy; go get him,'' then counted to three and pointed in the boy's direction (\href{https://supreme.justia.com/cases/federal/us/508/476/}{Justia Opinion}). The group ran across the street, severely beat the boy, and stole his tennis shoes. The victim was rendered unconscious and remained in a coma for four days (\href{https://www.courtlistener.com/opinion/112879/wisconsin-v-mitchell/}{CourtListener}).

Mitchell was charged with aggravated battery under Wis. Stat. \S\S~939.05 and 940.19(1m) (1989--1990), an offense that ordinarily carried a maximum sentence of two years' imprisonment under \S\S~940.19(1m) and 939.50(3)(e). Prosecutors also invoked Wisconsin's hate-crime penalty-enhancement statute, Wis. Stat. \S~939.645 (1989--1990), which increased the maximum penalty for an underlying felony by up to five years (and up to \$5,000 in fines) whenever the defendant ``[i]ntentionally selects the person against whom the crime \ldots\ is committed \ldots\ because of the race, religion, color, disability, sexual orientation, national origin or ancestry of that person'' (\href{https://supreme.justia.com/cases/federal/us/508/476/}{Justia Opinion}).

\section{Procedural History}
After a jury trial in the Circuit Court for Kenosha County, Mitchell was convicted of aggravated battery, and the jury returned a special verdict finding that he had intentionally selected his victim because of the boy's race (\href{https://supreme.justia.com/cases/federal/us/508/476/}{Justia Opinion}). Under \S~939.645(2)(c), this finding raised Mitchell's maximum possible prison term from two years to seven years. The Circuit Court sentenced Mitchell to four years' imprisonment and denied his motion for postconviction relief (\href{https://www.fire.org/supreme-court/wisconsin-v-todd-mitchell}{FIRE}).

Mitchell appealed, challenging \S~939.645 under the First Amendment as well as Fourteenth Amendment equal protection and vagueness grounds. The Wisconsin Court of Appeals affirmed, rejecting his First Amendment and vagueness challenges and holding his equal protection claim waived, 163 Wis. 2d 652, 473 N.W.2d 1 (1991) (\href{https://supreme.justia.com/cases/federal/us/508/476/}{Justia Opinion}).

The Wisconsin Supreme Court reversed, 169 Wis. 2d 153, 485 N.W.2d 807 (1992), with Justices Shirley S. Abrahamson and William A. Bablitch dissenting (\href{https://www.shirleyabrahamson.org/criminal-law/}{Shirley Abrahamson Archive}; \href{https://supreme.justia.com/cases/federal/us/508/476/}{Justia Opinion}). Relying on \textit{R.A.V. v. City of St. Paul}, 505 U.S. 377 (1992), the Wisconsin Supreme Court held that \S~939.645 violated the First Amendment directly by punishing ``offensive thought'' (the defendant's motive for selecting the victim) and was unconstitutionally overbroad because admitting prior speech to prove motive would chill protected expression; it declined to address the Fourteenth Amendment claims (\href{https://supreme.justia.com/cases/federal/us/508/476/}{Justia Opinion}). The U.S. Supreme Court granted certiorari (506 U.S. 1033 (1992)) to resolve a split among state high courts, heard oral argument on April 21, 1993, and issued its decision on June 11, 1993 (\href{https://www.oyez.org/cases/1992/92-515}{Oyez}).

\section{Judicial Votes}
\textbf{Vote Breakdown:} 9--0 (Unanimous) (\href{https://www.oyez.org/cases/1992/92-515}{Oyez}; \href{https://www.law.cornell.edu/supct/html/92-515.ZS.html}{Cornell LII}).
\begin{itemize}
    \item \textbf{Majority (9):} Chief Justice William H. Rehnquist (author), joined by Justices Byron R. White, Harry A. Blackmun, John Paul Stevens, Sandra Day O'Connor, Antonin Scalia, Anthony M. Kennedy, David H. Souter, and Clarence Thomas.
    \item \textbf{Concurrences (0):} None.
    \item \textbf{Dissents (0):} None.
\end{itemize}

\section{Holding}
\textbf{Judgment:} Reversed and remanded. The U.S. Supreme Court reversed the judgment of the Supreme Court of Wisconsin and remanded for further proceedings (\href{https://supreme.justia.com/cases/federal/us/508/476/}{Justia Opinion}).

\textbf{Rule Announced:} The First Amendment does not prohibit a state from enhancing the maximum criminal penalty for unprotected criminal conduct when the defendant intentionally selects the victim because of the victim's race or other protected status, nor does the First Amendment forbid the evidentiary use of a defendant's prior speech to prove motive or intent, subject to ordinary rules of evidence (\href{https://www.courtlistener.com/opinion/112879/wisconsin-v-mitchell/}{CourtListener}; \href{https://firstamendment.mtsu.edu/article/wisconsin-v-mitchell/}{First Amendment Encyclopedia}).

\textbf{Holding vs. Dicta:} The \textit{holding} is that applying Wis. Stat. \S~939.645 to enhance Mitchell's aggravated-battery sentence did not violate the First Amendment, either as an impermissible punishment of thought or under the overbreadth doctrine, and that the Court was not bound by the Wisconsin Supreme Court's characterization of the statute's practical First Amendment effect. The Court's observations about the societal harms of bias-motivated crimes (such as provoking retaliatory violence and community unrest) and its discussion of how minor misdemeanors like negligent motor-vehicle operation are unlikely to be racially motivated function as supporting reasoning and explanatory dicta; the Court also expressly declined to reach Mitchell's undeveloped Fourteenth Amendment equal protection and vagueness claims (\href{https://supreme.justia.com/cases/federal/us/508/476/}{Justia Opinion}).

\section{Analysis of Opinions}

\subsection{Majority Opinion (Chief Justice Rehnquist)}
Writing for a unanimous Court, Chief Justice Rehnquist dismantled each pillar of the Wisconsin Supreme Court's reasoning (\href{https://supreme.justia.com/cases/federal/us/508/476/}{Justia Opinion}):
\begin{enumerate}
    \item \textbf{Not Bound by State Court's Constitutional Characterization:} Although federal courts are bound by a state high court's construction of the literal meaning of a state statute (\textit{Terminiello v. Chicago}, 337 U.S. 1 (1949)), the Wisconsin Supreme Court had not defined a statutory word or phrase; rather, it characterized the statute's ``practical effect'' for First Amendment purposes. The Supreme Court remained free to form its own judgment on the law's operative First Amendment effect.
    \item \textbf{Conduct vs. Speech and the Role of Motive in Sentencing:} Rehnquist acknowledged that the State's claim that \S~939.645 punished only conduct did not end the inquiry, because the same battery received a harsher maximum penalty solely due to the defendant's discriminatory motive. Still, physical assault is not expressive conduct protected by the First Amendment (\textit{United States v. O'Brien}, 391 U.S. 367 (1968); \textit{NAACP v. Claiborne Hardware Co.}, 458 U.S. 886 (1982)). Sentencing judges and legislatures have long considered a defendant's motive---such as committing murder for pecuniary gain---to assess culpability. While \textit{Dawson v. Delaware}, 503 U.S. 159 (1992), forbids punishing a defendant's ``abstract beliefs'' unrelated to the crime, \textit{Barclay v. Florida}, 463 U.S. 939 (1983), confirms that racial animus directly tied to the offense may be considered even in capital sentencing.
    \item \textbf{Analogy to Antidiscrimination Laws and Distinction from \textit{R.A.V.}:} Motive plays the same role under \S~939.645 as under Title VII of the Civil Rights Act of 1964 (42 U.S.C. \S~2000e-2(a)(1)), which prohibits discriminatory conduct ``because of'' race or other protected traits and has repeatedly survived First Amendment challenges (\textit{Hishon v. King \& Spalding}, 467 U.S. 69 (1984)). Crucially, Rehnquist distinguished \textit{R.A.V. v. City of St. Paul}, 505 U.S. 377 (1992): the St. Paul ordinance was explicitly directed at \textit{expression} (prohibiting bias-motivated ``fighting words'' and messages), whereas Wisconsin's statute targeted \textit{unprotected criminal conduct}. Moreover, Wisconsin had an adequate, content-neutral justification beyond mere disagreement with bigoted beliefs: bias-inspired crimes are reasonably thought to inflict greater individual and societal harms, including distinct emotional trauma, retaliatory crimes, and community unrest (\href{https://supreme.justia.com/cases/federal/us/508/476/}{Justia Opinion}).
    \item \textbf{Rejection of Overbreadth and Evidentiary Speech Objections:} The Court dismissed the ``chilling effect'' theory as far too speculative, doubting that citizens would suppress bigoted views out of fear that those statements might later be offered at trial if they someday committed a violent felony. Furthermore, the First Amendment does not bar admitting a defendant's prior statements as evidence to prove an element of a crime, intent, or motive, subject to standard evidentiary safeguards for relevance and reliability (\textit{Haupt v. United States}, 330 U.S. 631 (1947); \textit{Price Waterhouse v. Hopkins}, 490 U.S. 228 (1989)) (\href{https://supreme.justia.com/cases/federal/us/508/476/}{Justia Opinion}).
\end{enumerate}

\subsection{Concurrences}
No Justice filed a separate concurring opinion (\href{https://supreme.justia.com/cases/federal/us/508/476/}{Justia Opinion}).

\subsection{Dissents}
No Justice filed a dissenting opinion (\href{https://supreme.justia.com/cases/federal/us/508/476/}{Justia Opinion}).

\section{Examples: Future Applications}
\begin{enumerate}
    \item \textbf{Same-Side (Constitutional under \textit{Mitchell}):} A federal hate-crime statute enhances the sentence for arson when the defendant intentionally burns down a house of worship because of the congregation's religion, and prosecutors introduce the defendant's prior social-media posts expressing hatred of that faith to prove motive. \textit{Analysis:} Like aggravated battery in \textit{Mitchell}, arson is unprotected criminal conduct, the enhancement targets bias-motivated victim/property selection inflicting heightened community harm, and prior speech is admissible as evidence of motive under \textit{Haupt} and \textit{Mitchell}.
    \item \textbf{Same-Side (Constitutional under \textit{Mitchell}):} A state enhances the criminal penalty for vandalism whenever a defendant intentionally targets property because of the owner's sexual orientation or disability. \textit{Analysis:} Vandalism is non-expressive criminal conduct outside First Amendment protection; increasing the statutory maximum based on discriminatory target selection mirrors Wis. Stat. \S~939.645(1)(b) and is squarely valid under \textit{Mitchell}.
    \item \textbf{Opposite-Side (Unconstitutional / Outside \textit{Mitchell}):} A city enacts an ordinance imposing a one-year jail sentence on anyone who displays a symbol or shouts epithets on a public sidewalk that are likely to arouse anger or alarm ``on the basis of race, religion, or gender,'' without requiring any underlying assault, trespass, or property damage. \textit{Analysis:} Because the law targets pure expression rather than an independently unlawful act of violence or property crime, it falls under \textit{R.A.V. v. City of St. Paul} as an unconstitutional content- and viewpoint-based restriction on speech, not \textit{Mitchell}.
    \item \textbf{Opposite-Side (Unconstitutional / Outside \textit{Mitchell}):} During sentencing for an ordinary, non-racially motivated burglary of a vacant warehouse, the prosecution introduces evidence that the defendant belongs to an extremist political group and holds racist views, solely to persuade the judge to impose a harsher sentence. \textit{Analysis:} Under \textit{Dawson v. Delaware}, expressly reaffirmed in \textit{Mitchell}, a sentencing court violates the First Amendment by punishing a defendant's abstract beliefs or associations when they bear no nexus to the commission of the underlying offense.
    \item \textbf{Fence-Sitter (Close / Contested Case):} A state applies a hate-crime penalty enhancement to the misdemeanor offense of ``disorderly conduct'' when a protester screams racial slurs at a counter-protester during a heated public rally and briefly shoves the person's sign. \textit{Analysis:} This sits directly on the contested fault line between \textit{R.A.V.} and \textit{Mitchell}. If the underlying conviction rests primarily on verbal epithets or expressive protest activity, enhancing the penalty for racial bias risks punishing protected expression in violation of \textit{R.A.V.}; if the conviction rests strictly on an unprotected physical battery and the slurs serve only as evidence of victim selection, \textit{Mitchell} would uphold the enhancement.
\end{enumerate}

\section{Critique}
\textbf{Originalist, Textualist, and Civil-Libertarian Critiques:} Free-speech formalists and civil-libertarian scholars---including Susan Gellman (who co-authored Mitchell's Supreme Court brief), Steven G. Gey, and commentators in the \textit{Vermont Law Review} and \textit{Northern Illinois University Law Review}---criticize \textit{Mitchell} for collapsing the constitutional boundary between punishing criminal acts and punishing disfavored thoughts (\href{https://lawreview.vermontlaw.edu/wp-content/uploads/2026/01/18VtLRev771.pdf}{Vermont Law Review}; \href{https://huskiecommons.lib.niu.edu/cgi/viewcontent.cgi?article=1569&context=niulr}{Northern Illinois University Law Review}). Critics argue that because the underlying battery is already fully punished, the \textit{only} factor triggering the extra five years of prison exposure is the defendant's ideological motive---an internal belief protected by the First Amendment (\href{https://lawecommons.luc.edu/cgi/viewcontent.cgi?article=1263&context=facpubs}{Loyola University Chicago Law}). From a textualist and doctrinal-consistency standpoint, critics also contend that \textit{Mitchell} sits in sharp tension with \textit{R.A.V. v. City of St. Paul}: both laws single out bias-motivated bigotry for extra punishment, yet \textit{Mitchell} allowed the legislature to evade strict scrutiny merely by attaching the bias penalty to an existing criminal statute rather than drafting a standalone speech ordinance (\href{https://lawreview.vermontlaw.edu/wp-content/uploads/2026/01/18VtLRev771.pdf}{Vermont Law Review}). Some scholars (such as James Jacobs and Kimberly Potter) further question the empirical claim that hate-crime enhancements reliably deter bias violence without deepening social divisions (\href{https://lawecommons.luc.edu/cgi/viewcontent.cgi?article=1263&context=facpubs}{Loyola University Chicago Law}).

\textbf{Progressive and Doctrinal Perspectives:} Conversely, progressive and civil-rights scholars such as Frederick M. Lawrence, alongside major civil-rights organizations including the ACLU, Anti-Defamation League, and Lawyers' Committee for Civil Rights Under Law, praise \textit{Mitchell} as a pragmatic and doctrinally sound resolution (\href{https://lawecommons.luc.edu/cgi/viewcontent.cgi?article=1263&context=facpubs}{Loyola University Chicago Law}; \href{https://scholarlycommons.law.wlu.edu/cgi/viewcontent.cgi?article=1015&context=crsj}{Washington and Lee Civil Rights and Social Justice Journal}). They emphasize that bias-motivated violence is qualitatively more destructive than parallel ordinary crimes because it terrorizes an entire targeted community, undermines equal citizenship, and risks cycles of retaliatory violence. By grounding its decision in traditional sentencing doctrine and Title VII antidiscrimination principles, the unanimous Court gave legislatures a clear, workable blueprint for combating hate violence without eroding First Amendment protection for non-violent offensive speech (\href{https://firstamendment.mtsu.edu/article/wisconsin-v-mitchell/}{First Amendment Encyclopedia}).

\textbf{Underlying Competing Values:} \textit{Mitchell} balances \textit{freedom of thought and viewpoint neutrality} (preventing the state from using criminal punishment to enforce ideological orthodoxy or chill unpopular speech) against \textit{equal protection, public safety, and community dignity} (empowering legislatures to deter and punish discriminatory violence that inflicts heightened harm on both individual victims and marginalized groups).

\section{Key Quotations}

\begin{quote}
``Thus, a physical assault is not by any stretch of the imagination expressive conduct protected by the First Amendment.''
\end{quote}
\noindent --- Chief Justice William H. Rehnquist, Majority Opinion, 508 U.S. at 484 (\href{https://supreme.justia.com/cases/federal/us/508/476/}{Justia}; \href{https://www.courtlistener.com/opinion/112879/wisconsin-v-mitchell/}{CourtListener}).

\begin{quote}
``But it is equally true that a defendant's abstract beliefs, however obnoxious to most people, may not be taken into consideration by a sentencing judge.''
\end{quote}
\noindent --- Chief Justice William H. Rehnquist, Majority Opinion, 508 U.S. at 485--486 (\href{https://supreme.justia.com/cases/federal/us/508/476/}{Justia}; \href{https://www.courtlistener.com/opinion/112879/wisconsin-v-mitchell/}{CourtListener}).

\begin{quote}
``But whereas the ordinance struck down in R. A. V. was explicitly directed at expression (i. e., `speech' or `messages'), id., at 392, the statute in this case is aimed at conduct unprotected by the First Amendment.''
\end{quote}
\noindent --- Chief Justice William H. Rehnquist, Majority Opinion, 508 U.S. at 487 (\href{https://supreme.justia.com/cases/federal/us/508/476/}{Justia}; \href{https://www.courtlistener.com/opinion/112879/wisconsin-v-mitchell/}{CourtListener}).

\begin{quote}
``Moreover, the Wisconsin statute singles out for enhancement bias-inspired conduct because this conduct is thought to inflict greater individual and societal harm.''
\end{quote}
\noindent --- Chief Justice William H. Rehnquist, Majority Opinion, 508 U.S. at 487--488 (\href{https://supreme.justia.com/cases/federal/us/508/476/}{Justia}; \href{https://www.courtlistener.com/opinion/112879/wisconsin-v-mitchell/}{CourtListener}).

\begin{quote}
``The First Amendment, moreover, does not prohibit the evidentiary use of speech to establish the elements of a crime or to prove motive or intent.''
\end{quote}
\noindent --- Chief Justice William H. Rehnquist, Majority Opinion, 508 U.S. at 489 (\href{https://supreme.justia.com/cases/federal/us/508/476/}{Justia}; \href{https://www.courtlistener.com/opinion/112879/wisconsin-v-mitchell/}{CourtListener}).

\end{document}
